\section{Proofs for the Debiased Estimator}\label{app:estimator}

Throughout this appendix, condition on one frozen target and write
$g=\nabla J(\theta)$.  The complete trajectory summaries $\mathsf X_1,\mathsf X_2,\ldots$ are independent and identically distributed under this conditional law.
For a realization $x=(V^+,V^-,G_x)$, write $I_{x,\sigma}(y)=\1\{V^\sigma>y\}$.

\subsection{Second-order expansion and full/half coupling}

\begin{lemma}[Coupled full/half cancellation]\label{lem:level-cancellation}
Under Assumptions \ref{ass:trajectory}--\ref{ass:cpt-regularity}, fix one frozen target and parameter and write $g:=\nabla J(\theta)$.  There exist a measurable $\mathfrak h$ and constants $C_{\rm rem},C_U<\infty$ such that
\begin{equation}\label{eq:loo-expansion}
 U_n^{\rm raw}-g
 =\frac1n\sum_{i=1}^{n}\mathfrak h(\mathsf X_i)+\mathcal B_n,
 \qquad
 \E\mathfrak h(\mathsf X_i)=0,
 \qquad
 \E\norm{\mathcal B_n}_2^2\le\frac{C_{\rm rem}}{n^2},
\end{equation}
and consequently
\begin{equation}\label{eq:level-second-moment}
 \E\norm{D_{n,\ell}^{\rm raw}}_2^2
 \le\frac{9C_{\rm rem}}{n^2}2^{-2\ell},
 \qquad \ell\ge1,
\end{equation}
while
\begin{equation}\label{eq:centered-loo-variance}
 \tr\Cov(U_n^{\rm cen})\le\frac{C_U}{n}.
\end{equation}
\end{lemma}

\begin{proof}[Proof of \Cref{lem:level-cancellation}]
To prove \eqref{eq:loo-expansion}, it is enough to isolate the first-order Hoeffding projection and show that both the degenerate two-sample term and the Taylor remainder are $O_{L^2}(n^{-1})$.  For a side $\sigma$ and threshold $y$, write
\[
 S_\sigma(y)=\Pp(V^\sigma>y),
 \qquad
 e_{-i,\sigma}(y)=\widehat S_{-i,n}^\sigma(y)-S_\sigma(y).
\]
Taylor's theorem applied to $(w_\sigma^\eps)'$ gives
\begin{equation}\label{eq:weight-taylor}
 (w_\sigma^\eps)'(S_\sigma+e)
 =(w_\sigma^\eps)'(S_\sigma)
 +(w_\sigma^\eps)''(S_\sigma)e
 +\operatorname{Rem}_\sigma(e),
 \qquad
 \abs{\operatorname{Rem}_\sigma(e)}\le\frac{C_3}{2}e^2.
\end{equation}
Substituting \eqref{eq:weight-taylor} into \eqref{eq:raw-loo} splits the statistic into a first-order trajectory term, an ordered two-sample term, and a Taylor remainder.

Define
\begin{align*}
 \mathfrak h_0(x)
 &=G_x\sum_\sigma \varsigma_\sigma
 \int_0^{B_v}(w_\sigma^\eps)'(S_\sigma(y))I_{x,\sigma}(y)\,dy,\\
 \mathcal K(x,x')
 &=G_x\sum_\sigma \varsigma_\sigma
 \int_0^{B_v}(w_\sigma^\eps)''(S_\sigma(y))
 I_{x,\sigma}(y)\{I_{x',\sigma}(y)-S_\sigma(y)\}\,dy.
\end{align*}
By \eqref{eq:score-gradient-identity} and $\E G=0$, $\E \mathfrak h_0(\mathsf X)=g$.  Moreover
$\E\{\mathcal K(\mathsf X,\mathsf X')\mid \mathsf X\}=0$ for independent $\mathsf X,\mathsf X'$, because
$\E\{I_{\mathsf X',\sigma}(y)-S_\sigma(y)\}=0$.  Let
\[
 \mathfrak h_1(x')=\E_{\mathsf X}\{\mathcal K(\mathsf X,x')\}.
\]
Then $\E \mathfrak h_1(\mathsf X')=0$.  Put
$\mathcal K^\circ(x,x')=\mathcal K(x,x')-\mathfrak h_1(x')$.  Both first-order projections of $\mathcal K^\circ$ vanish:
\[
 \E\{\mathcal K^\circ(\mathsf X,\mathsf X')\mid \mathsf X\}=0,
 \qquad
 \E\{\mathcal K^\circ(\mathsf X,\mathsf X')\mid \mathsf X'\}=0.
\]
The ordered two-sample term therefore has the exact decomposition
\begin{equation}\label{eq:ordered-hoeffding}
 \frac{1}{n(n-1)}\sum_{i\ne j}\mathcal K(\mathsf X_i,\mathsf X_j)
 =\frac1n\sum_{j=1}^{n}\mathfrak h_1(\mathsf X_j)
 +\mathcal D_n,
\end{equation}
where
\[
 \mathcal D_n=\frac{1}{n(n-1)}\sum_{i\ne j}\mathcal K^\circ(\mathsf X_i,\mathsf X_j).
\]
The kernel is square-integrable because the indicators are bounded, the integration interval is finite, $(w_\sigma^\eps)''$ is bounded, and $\E\norm{G}_2^2<\infty$.  Degeneracy implies that covariances between two ordered kernel terms vanish unless their ordered index pairs coincide or are reversals of one another.  Counting these $O(n^2)$ contributing pairs against the normalization $n^2(n-1)^2$ gives
\begin{equation}\label{eq:degenerate-remainder}
 \E\norm{\mathcal D_n}_2^2\le\frac{C_D}{n^2}
\end{equation}
for a finite constant $C_D$.

We next bound the Taylor remainder.  For each $i$, $e_{-i,\sigma}$ is an empirical mean based on $n-1$ Bernoulli variables independent of $\mathsf X_i$.  Uniformly in $y$,
\begin{equation}\label{eq:bernoulli-fourth-moment}
 \E\abs{e_{-i,\sigma}(y)}^4\le\frac{C_e}{(n-1)^2}
\end{equation}
for a universal finite $C_e$.  By \eqref{eq:weight-taylor}, Cauchy--Schwarz in the $y$-integral, independence of $G_i$ from the leave-one-out empirical survival, and $\E\norm{G_i}_2^2<\infty$, the $i$th Taylor contribution has second moment at most $C_T/n^2$.  Jensen's inequality for the average over $i$ then gives a second-moment bound of the same order for the averaged Taylor remainder $\mathcal T_n$:
\begin{equation}\label{eq:taylor-remainder-bound}
 \E\norm{\mathcal T_n}_2^2\le\frac{C_T}{n^2}.
\end{equation}

Combining the first-order term, \eqref{eq:ordered-hoeffding}, and the Taylor remainder yields
\[
 U_n^{\rm raw}-g
 =\frac1n\sum_{i=1}^{n}\mathfrak h(\mathsf X_i)+\mathcal B_n,
 \qquad
 \mathfrak h=\mathfrak h_0-g+\mathfrak h_1,
 \qquad
 \mathcal B_n=\mathcal D_n+\mathcal T_n.
\]
The vector function $\mathfrak h$ has mean zero, and \eqref{eq:degenerate-remainder}--\eqref{eq:taylor-remainder-bound} imply
$\E\norm{\mathcal B_n}_2^2\le C_{\rm rem}/n^2$ after increasing the finite constant $C_{\rm rem}$.  This proves \eqref{eq:loo-expansion}.

Now take $n_\ell=n2^\ell$ and evaluate the full-sample statistic and its two half-sample statistics on the same $n_\ell$ trajectories.  Applying \eqref{eq:loo-expansion} to the full sample and to each half gives
\[
 U_{n_\ell}^{\rm raw}
 =g+\frac1{n_\ell}\sum_{i=1}^{n_\ell}\mathfrak h(\mathsf X_i)+\mathcal B_{n_\ell},
\]
and
\[
 \frac12\{U_{n_\ell/2}^{\rm raw,(1)}+U_{n_\ell/2}^{\rm raw,(2)}\}
 =g+\frac1{n_\ell}\sum_{i=1}^{n_\ell}\mathfrak h(\mathsf X_i)
 +\frac12\{\mathcal B_{n_\ell/2}^{(1)}+\mathcal B_{n_\ell/2}^{(2)}\}.
\]
The influence terms cancel exactly.  Hence
\[
 D_{n,\ell}^{\rm raw}
 =\mathcal B_{n_\ell}-\frac12\{\mathcal B_{n_\ell/2}^{(1)}+\mathcal B_{n_\ell/2}^{(2)}\}.
\]
Minkowski's inequality in $L^2$ gives
\begin{align*}
 (\E\norm{D_{n,\ell}^{\rm raw}}_2^2)^{1/2}
 &\le\frac{\sqrt{C_{\rm rem}}}{n_\ell}
 +\frac12\frac{2\sqrt{C_{\rm rem}}}{n_\ell}
 +\frac12\frac{2\sqrt{C_{\rm rem}}}{n_\ell}
 =\frac{3\sqrt{C_{\rm rem}}}{n_\ell}.
\end{align*}
Squaring and substituting $n_\ell=n2^\ell$ proves \eqref{eq:level-second-moment}.

It remains to prove \eqref{eq:centered-loo-variance}.  View $U_n^{\rm cen}$ as a symmetric function of $\mathsf X_1,\ldots,\mathsf X_n$, and let $U_n^{\rm cen,(i)}$ be the statistic after replacing $\mathsf X_i$ by an independent copy $\mathsf X_i'$.  If $G_i'$ denotes the score component of $\mathsf X_i'$, the coefficient in one centered LOO summand is bounded in absolute value by $2B_vC_1$, so the contribution from the replaced summand has norm at most
\[
 \frac{2B_vC_1}{n}(\norm{G_i}_2+\norm{G_i'}_2).
\]
For $j\ne i$, replacing $\mathsf X_i$ changes each leave-one-out survival entering the $j$th summand by at most $(n-1)^{-1}$ pointwise.  Since
\[
 \left|\frac{\partial}{\partial p}\{(w_\sigma^\eps)'(p)(I-p)\}\right|
 \le C_1+C_2,
\]
the scalar coefficient multiplying $G_j$ changes by at most
$2B_v(C_1+C_2)/(n-1)$.  Hence, for a deterministic constant $C_*$ depending only on $B_v,C_1,C_2$,
\begin{align*}
 \norm{U_n^{\rm cen}-U_n^{\rm cen,(i)}}_2
 &\le \frac{C_*}{n}(\norm{G_i}_2+\norm{G_i'}_2)
 +\frac{C_*}{n(n-1)}\sum_{j\ne i}\norm{G_j}_2.
\end{align*}
Using $(\sum_{j\ne i}x_j)^2\le(n-1)\sum_{j\ne i}x_j^2$ and the uniform second-moment bound for the score, there is a finite constant $C_E$ such that
\[
 \E\norm{U_n^{\rm cen}-U_n^{\rm cen,(i)}}_2^2
 \le\frac{C_E}{n^2}
\]
uniformly in $i$.  The vector Efron--Stein inequality yields
\[
 \tr\Cov(U_n^{\rm cen})
 =\E\norm{U_n^{\rm cen}-\E U_n^{\rm cen}}_2^2
 \le\frac12\sum_{i=1}^{n}
 \E\norm{U_n^{\rm cen}-U_n^{\rm cen,(i)}}_2^2
 \le\frac{C_U}{n}
\]
with $C_U=C_E/2$.  This completes the proof.
\end{proof}


\subsection{Centering and finite-sample bias}

\begin{lemma}[Centering and finite-sample bias]\label{lem:centering-expectation}
For $n\ge2$,
\begin{equation}\label{eq:centered-raw-mean}
 \E U_n^{\rm cen}=\E U_n^{\rm raw}=: \bar U_n,
\end{equation}
and
\begin{equation}\label{eq:loo-bias-order}
 \norm{\bar U_n-g}_2\le \frac{\sqrt{C_{\rm rem}}}{n},
 \qquad g:=\nabla J(\theta).
\end{equation}
\end{lemma}

\begin{proof}[Proof of \Cref{lem:centering-expectation}]
To compare the two LOO statistics, subtract \eqref{eq:centered-loo} from \eqref{eq:raw-loo}.  The difference is
\begin{equation}\label{eq:centering-control}
 U_n^{\rm raw}-U_n^{\rm cen}
 =\frac1n\sum_{i=1}^{n}G_i H_{-i}^{\rm ctr},
\end{equation}
where
\[
 H_{-i}^{\rm ctr}
 =\sum_{\sigma\in\{+,-\}}\varsigma_\sigma
 \int_0^{B_v}
 (w_\sigma^\eps)'(\widehat S_{-i,n}^\sigma(y))
 \widehat S_{-i,n}^\sigma(y)\,dy.
\]
The coefficient $H_{-i}^{\rm ctr}$ depends only on trajectories with indices different from $i$.  It is therefore independent of $G_i$.  The score identity gives $\E G_i=0$, while Assumptions \ref{ass:trajectory}--\ref{ass:cpt-regularity} make $G_iH_{-i}^{\rm ctr}$ integrable.  Hence every summand in \eqref{eq:centering-control} has mean zero and
$\E U_n^{\rm cen}=\E U_n^{\rm raw}=\bar U_n$.

It remains to control $\bar U_n-g$.  This follows from the expansion in \Cref{lem:level-cancellation}.  Taking expectations in \eqref{eq:loo-expansion} and using $\E\mathfrak h(\mathsf X)=0$ yields
\[
 \norm{\bar U_n-g}_2
 =\norm{\E \mathcal B_n}_2
 \le(\E\norm{\mathcal B_n}_2^2)^{1/2}
 \le\frac{\sqrt{C_{\rm rem}}}{n}.
\]
This proves \eqref{eq:loo-bias-order}.
\end{proof}


\subsection{Unbiasedness, variance, cost, and the CLT}

\begin{proof}[Proof of \Cref{thm:estimator}]
We first prove exact unbiasedness.  By \Cref{lem:centering-expectation}, the centered base has mean $\bar U_n$.  For the raw correction,
\begin{equation}\label{eq:delta-mean}
 \E D_{n,\ell}^{\rm raw}
 =\bar U_{n2^\ell}-\bar U_{n2^{\ell-1}}.
\end{equation}
The level bound \eqref{eq:level-second-moment} implies
\[
 \sum_{\ell\ge1}\E\norm{D_{n,\ell}^{\rm raw}}_2
 \le\frac{3\sqrt{C_{\rm rem}}}{n}\sum_{\ell\ge1}2^{-\ell}<\infty,
\]
so the expectation can be interchanged with the infinite randomized series.  Since \eqref{eq:loo-bias-order} gives $\bar U_{n2^{\ell_{\max}}}\to g$,
\begin{align*}
 \E Z_n^{\db}
 &=\bar U_n+\sum_{\ell\ge1}\E D_{n,\ell}^{\rm raw}\\
 &=\bar U_n+\lim_{\ell_{\max}\to\infty}(\bar U_{n2^{\ell_{\max}}}-\bar U_n)
 =g.
\end{align*}
This proves \eqref{eq:estimator-unbiased}.

For the second moment of the correction, independence of the activation and level variables gives
\begin{align*}
 \E\left\|\frac{\mathsf A}{p_{\rm act}\nu_{\ell^\star}}D_{n,\ell^\star}^{\rm raw}\right\|_2^2
 &=\frac{1}{p_{\rm act}}\sum_{\ell\ge1}
 \frac{\E\norm{D_{n,\ell}^{\rm raw}}_2^2}{\nu_\ell}\\
 &\le\frac{9C_{\rm rem}}{p_{\rm act} n^2}
 \sum_{\ell\ge1}\frac{4^{-\ell}}{\nu_\ell}
 =\frac{9C_{\rm rem}\Xi_b}{p_{\rm act} n^2}.
\end{align*}
The series $\Xi_b$ is finite because $b<2$.  The base sample and correction sample are independent, so their covariance is zero.  Combining the correction bound with \eqref{eq:centered-loo-variance} gives \eqref{eq:single-output-variance}.

One base output costs exactly $n$ complete trajectories.  Conditional on activation and level $\ell$, the correction uses one fresh batch of $n2^\ell$ trajectories; the two halves reuse this same batch.  Hence
\begin{align*}
 \E\operatorname{cost}(Z_n^{\db})
 &=n+p_{\rm act} n\sum_{\ell\ge1}\nu_\ell2^\ell
 =n(1+p_{\rm act} \Upsilon_b),
\end{align*}
where $\Upsilon_b<\infty$ because $b>1$.  Direct summation of the geometric series gives the closed forms in \eqref{eq:level-series-constants}.

Finally, if $Z_{n,1}^{\db},\ldots,Z_{n,M}^{\db}$ are independent copies, exact unbiasedness and independence imply
\[
 \E\norm{\widehat g_{M,n}^{\db}-g}_2^2
 =\frac1M\tr\Cov(Z_n^{\db})
 \le\frac{C_{\rm var}(n,p_{\rm act},b)}{M}.
\]
This is \eqref{eq:estimator-mse}.
\end{proof}

\begin{proof}[Proof of \Cref{thm:clt}]
For fixed target and fixed $(n,p_{\rm act},b)$, \Cref{thm:estimator} gives iid vectors $Z_{n,m}^{\db}$ with mean $g$ and finite covariance $\Sigma_n$.  For any fixed $t\in\R^d$, the scalar variables
$t^\top(Z_{n,m}^{\db}-g)$ are iid with finite variance.  The classical scalar central limit theorem therefore gives
\[
 \sqrt M\,t^\top(\widehat g_{M,n}^{\db}-g)
 \rightsquigarrow \mathcal N(0,t^\top\Sigma_nt).
\]
The Cram\'er--Wold device proves \eqref{eq:clt}.

Because $\E\norm{Z_n^{\db}}_2^2<\infty$, every entry of $Z_n^{\db}(Z_n^{\db})^\top$ is integrable.  The law of large numbers therefore implies
$\widehat\Sigma_{M,n}\to\Sigma_n$ in probability entrywise, hence in operator norm in fixed dimension.  If $\Sigma_n$ is nonsingular, operator-norm consistency implies
$\Pp\{\lambda_{\min}(\widehat\Sigma_{M,n})>\lambda_{\min}(\Sigma_n)/2\}\to1$.
On this event the inverse square-root map is well defined and continuous in a neighborhood of $\Sigma_n$.  Slutsky's theorem applied to \eqref{eq:clt} then proves \eqref{eq:studentized-clt}.
\end{proof}
